%%%
%%% Radical "⽹"
%%%
\section*{Radical 122: ``⽹'' (⺲、罓、⺳)}\addcontentsline{toc}{section}{Radical 122: ⽹,⺲、罓、⺳}\addcontentsline{loh}{figure}{\#\#\#\# 122: ⽹}

%%%%%%%%%% 网 %%%%%%%%%%
\subsection*{网}\addcontentsline{loh}{figure}{网}

\begin{Entry}{网}{6}{⽹}[Kangxi 122]
  \begin{Phonetics}{网}{wang3}[][HSK 2]
    \definition[张]{s.}{rede; um dispositivo feito de corda ou barbante para capturar peixes e pássaros | algo que parece uma rede | rede; uma rede de organizações; um sistema}
    \definition{v.}{pegar com uma rede | cobrir como com uma rede}
  \end{Phonetics}
\end{Entry}

\begin{Entry}{网上}{6,3}{⽹,⼀}
  \begin{Phonetics}{网上}{wang3shang4}[][HSK 1]
    \definition{s.}{\emph{on"-line}; refere"-se a acessar a \emph{Internet} através de um computador ou celular para pesquisar e consultar informações na rede}
  \end{Phonetics}
\end{Entry}

\begin{Entry}{网上银行}{6,3,11,6}{⽹,⼀,⾦,⾏}
  \begin{Phonetics}{网上银行}{wang3shang4yin2hang2}
    \definition[个]{s.}{\emph{e"-bank}; banco \emph{on"-line} | acesso a operações bancárias via \emph{Internet}}
  \seealsoref{网银}{wang3yin2}
  \end{Phonetics}
\end{Entry}

\begin{Entry}{网友}{6,4}{⽹,⼜}
  \begin{Phonetics}{网友}{wang3you3}[][HSK 1]
    \definition[位,个,名]{s.}{internauta; usuário da \emph{Internet}; amigos que se conhecem pela Internet; também usado como forma de tratamento entre internautas}
  \end{Phonetics}
\end{Entry}

\begin{Entry}{网民}{6,5}{⽹,⽒}
  \begin{Phonetics}{网民}{wang3min2}[][HSK 7-9]
    \definition[位,名]{s.}{internauta; geralmente se refere a um usuário de rede de computadores}
  \end{Phonetics}
\end{Entry}

\begin{Entry}{网页}{6,6}{⽹,⾴}
  \begin{Phonetics}{网页}{wang3ye4}[][HSK 6]
    \definition[个]{s.}{site; página da web; \emph{website}; \emph{web page}}
  \end{Phonetics}
\end{Entry}

\begin{Entry}{网吧}{6,7}{⽹,⼝}
  \begin{Phonetics}{网吧}{wang3ba1}[][HSK 6]
    \definition[家,间]{s.}{cybercafé; \emph{Internet} café; refere"-se a um local comercial aberto ao público que utiliza redes de computadores para fornecer serviços de navegação, consulta e outras informações}
  \end{Phonetics}
\end{Entry}

\begin{Entry}{网址}{6,7}{⽹,⼟}
  \begin{Phonetics}{网址}{wang3zhi3}[][HSK 4]
    \definition[个]{s.}{\emph{website}; endereço da \emph{web}; endereço de um \emph{site} na \emph{Internet}, que os usuários podem acessar, consultar e obter recursos de informações nesse \emph{site} clicando nele}
  \end{Phonetics}
\end{Entry}

\begin{Entry}{网际网络}{6,7,6,9}{⽹,⾩,⽹,⽷}
  \begin{Phonetics}{网际网络}{wang3ji4 wang3luo4}
    \definition*{s.}{Internet}
  \seealsoref{互联网}{hu4lian2wang3}
  \seealsoref{网际网路}{wang3ji4wang3lu4}
  \seealsoref{网路}{wang3lu4}
  \end{Phonetics}
\end{Entry}

\begin{Entry}{网际网路}{6,7,6,13}{⽹,⾩,⽹,⾜}
  \begin{Phonetics}{网际网路}{wang3ji4wang3lu4}
    \definition*{s.}{Internet}
  \seealsoref{互联网}{hu4lian2wang3}
  \seealsoref{网际网络}{wang3ji4 wang3luo4}
  \seealsoref{网路}{wang3lu4}
  \end{Phonetics}
\end{Entry}

\begin{Entry}{网点}{6,9}{⽹,⽕}
  \begin{Phonetics}{网点}{wang3dian3}[][HSK 7-9]
    \definition{s.}{pontos de venda; rede (de estabelecimento comercial)}
  \end{Phonetics}
\end{Entry}

\begin{Entry}{网络}{6,9}{⽹,⽷}
  \begin{Phonetics}{网络}{wang3luo4}[][HSK 4]
    \definition{s.}{rede; um sistema que consiste em ramificações interconectadas; em um sistema elétrico, um circuito ou parte de um circuito que consiste em vários elementos que permitem a transmissão de sinais elétricos de acordo com determinados requisitos | rede; rede de computadores}
  \end{Phonetics}
\end{Entry}

\begin{Entry}{网站}{6,10}{⽹,⽴}
  \begin{Phonetics}{网站}{wang3zhan4}[][HSK 2]
    \definition[个,家]{s.}{\emph{web}; \emph{website}; um site virtual na Internet para uma organização ou indivíduo, geralmente consistindo em uma página inicial e muitas páginas da web}
  \end{Phonetics}
\end{Entry}

\begin{Entry}{网罟}{6,10}{⽹,⽹}
  \begin{Phonetics}{网罟}{wang3gu3}
    \definition{s.}{Figurativo:  a rede da justiça (法网) | rede usada para pegar peixes (ou outros animais, como pássaros)}
  \seealsoref{法网}{fa3wang3}
  \end{Phonetics}
\end{Entry}

\begin{Entry}{网球}{6,11}{⽹,⽟}
  \begin{Phonetics}{网球}{wang3qiu2}[][HSK 2]
    \definition[个,颗,些]{s.}{tênis (esporte) | bola de tênis}
  \end{Phonetics}
\end{Entry}

\begin{Entry}{网银}{6,11}{⽹,⾦}
  \begin{Phonetics}{网银}{wang3yin2}
    \definition{s.}{\emph{e"-bank}; banco eletrônico; banco pela \emph{Internet} | banco online; abreviação de 网上银行}
  \seealsoref{网上银行}{wang3shang4yin2hang2}
  \end{Phonetics}
\end{Entry}

\begin{Entry}{网路}{6,13}{⽹,⾜}
  \begin{Phonetics}{网路}{wang3lu4}
    \definition*{s.}{Internet}
  \seealsoref{互联网}{hu4lian2wang3}
  \seealsoref{网际网络}{wang3ji4 wang3luo4}
  \seealsoref{网际网路}{wang3ji4wang3lu4}
  \end{Phonetics}
\end{Entry}

%%%%%%%%%% 罕 %%%%%%%%%%
\subsection*{罕}\addcontentsline{loh}{figure}{罕}

\begin{Entry}{罕}{7}{⽹}
  \begin{Phonetics}{罕}{han3}
    \definition*{s.}{Sobrenome: Han}
    \definition{adj.}{raro; escasso | raro; incomum}
  \end{Phonetics}
\end{Entry}

\begin{Entry}{罕见}{7,4}{⽹,⾒}
  \begin{Phonetics}{罕见}{han3jian4}[][HSK 7-9]
    \definition{adj.}{raro; raramente visto}
  \end{Phonetics}
\end{Entry}

%%%%%%%%%% 罗 %%%%%%%%%%
\subsection*{罗}\addcontentsline{loh}{figure}{罗}

\begin{Entry}{罗}{8}{⽹}
  \begin{Phonetics}{罗}{luo2}[][HSK 7-9]
    \definition*{s.}{Sobrenome: Luo}
    \definition{clas.}{uma grosa; uma bruta; doze dúzias; 144 unidades}
    \definition{s.}{uma rede para capturar pássaros; redes de proteção contra pássaros | peneira; coador; tela | uma espécie de gaze de seda; tecidos de seda com textura solta}
    \definition{v.}{capturar pássaros com uma rede; prender pássaros em redes | espalhar; exibir; mostrar | coletar; reunir; recrutar; invocar | peneirar}
  \end{Phonetics}
\end{Entry}

%%%%%%%%%% 罚 %%%%%%%%%%
\subsection*{罚}\addcontentsline{loh}{figure}{罚}

\begin{Entry}{罚}{9}{⽹}
  \begin{Phonetics}{罚}{fa2}[][HSK 5]
    \definition{s.}{punição; penalidade}
    \definition{v.}{punir; penalizar; multar; confiscar}
  \end{Phonetics}
\end{Entry}

\begin{Entry}{罚款}{9,12}{⽹,⽋}
  \begin{Phonetics}{罚款}{fa2/kuan3}[][HSK 5]
    \definition[笔,次,宗]{s.}{multa; penalidade; refere"-se ao dinheiro pago por uma pessoa ou entidade de acordo com as disposições de um delito ou violação de contrato ou contrato}
    \definition{v.+compl.}{multar; penalizar; exigir, de acordo com os regulamentos, uma determinada quantia de dinheiro de uma pessoa ou entidade que tenha violado a lei ou descumprido um regulamento ou contrato}
  \end{Phonetics}
\end{Entry}

%%%%%%%%%% 罢 %%%%%%%%%%
\subsection*{罢}\addcontentsline{loh}{figure}{罢}

\begin{Entry}{罢}{10}{⽹}
  \begin{Phonetics}{罢}{ba4}
    \definition{v.}{parar; cessar | revogar; destituir; encerrar | terminar | abandonar uma ideia; esqueçer sobre algo; deixar estar (passar)}
  \end{Phonetics}
  \begin{Phonetics}{罢}{ba5}
    \definition{part.}{partícula final, a mesma que 吧}
  \seealsoref{吧}{ba5}
  \end{Phonetics}
\end{Entry}

\begin{Entry}{罢了}{10,2}{⽹,⼅}
  \begin{Phonetics}{罢了}{ba4le5}[][HSK 6]
    \definition{part.}{usado no final de uma frase, significa 仅此而已, geralmente seguido de 无非, 不过, 只是}
  \seealsoref{不过}{bu2guo4}
  \seealsoref{仅此而已}{jin3ci3'er2yi3}
  \seealsoref{无非}{wu2fei1}
  \seealsoref{只是}{zhi3shi4}
  \synonymref{而已}{er2yi3}
  \synonymref{算了}{suan4le5}
  \antonymref{继续}{ji4xu4}
  \antonymref{坚持}{jian1chi2}
  \end{Phonetics}
  \begin{Phonetics}{罢了}{ba4liao3}
    \definition{part.}{uma partícula modal indicando (não se preocupe, ok)}
  \end{Phonetics}
\end{Entry}

\begin{Entry}{罢工}{10,3}{⽹,⼯}
  \begin{Phonetics}{罢工}{ba4gong1}[][HSK 6]
    \definition{v.}{parar de trabalhar; entrar em greve; abandonar o emprego}
  \synonymref{停工}{ting2gong1}
  \end{Phonetics}
\end{Entry}

\begin{Entry}{罢休}{10,6}{⽹,⼈}
  \begin{Phonetics}{罢休}{ba4xiu1}[][HSK 7-9]
    \definition{v.}{parar; desistir; deixar o assunto de lado; parar de fazer algo, enfatizando a determinação de parar de fazê"-lo}
  \synonymref{放弃}{fang4qi4}
  \synonymref{放任}{fang4ren4}
  \synonymref{结束}{jie2shu4}
  \synonymref{截止}{jie2zhi3}
  \synonymref{停止}{ting2zhi3}
  \antonymref{继续}{ji4xu4}
  \end{Phonetics}
\end{Entry}

\begin{Entry}{罢免}{10,7}{⽹,⼉}
  \begin{Phonetics}{罢免}{ba4mian3}[][HSK 7-9]
    \definition{v.}{destituir; remover do cargo; demitir alguém do seu posto | destituir do cargo uma pessoa eleita pelo eleitorado ou por um órgão representativo}
  \synonymref{解除}{jie3chu2}
  \synonymref{解雇}{jie3gu4}
  \synonymref{免职}{mian3/zhi2}
  \synonymref{免除}{mian3chu2}
  \antonymref{选举}{xuan3ju3}
  \end{Phonetics}
\end{Entry}

%%%%%%%%%% 罩 %%%%%%%%%%
\subsection*{罩}\addcontentsline{loh}{figure}{罩}

\begin{Entry}{罩}{13}{⽹}
  \begin{Phonetics}{罩}{zhao4}[][HSK 7-9]
    \definition{s.}{cobertura; sombra; capuz; invólucro | armadilha de bambu para peixes | gaiola; galinheiro; galinheiro de bambu; gaiolas menores para criação de galinhas | macacão; casaco de trabalho; macacão de trabalho | abrigo; alojamento | casaco exterior; sobretudo; macacão; roupa exterior; bata}
    \definition{v.}{cobrir; espalhar; embrulhar | disseminar; cobrir; prender; colocar por fora}
  \seealsoref{罩儿}{zhao4r5}
  \end{Phonetics}
\end{Entry}

\begin{Entry}{罩儿}{13,2}{⽹,⼉}
  \begin{Phonetics}{罩儿}{zhao4r5}
    \definition{s.}{cobertura; sombra; capuz; invólucro | macacão; casaco de trabalho; macacão de trabalho}
  \end{Phonetics}
\end{Entry}

%%%%%%%%%% 罪 %%%%%%%%%%
\subsection*{罪}\addcontentsline{loh}{figure}{罪}

\begin{Entry}{罪}{13}{⽹}
  \begin{Phonetics}{罪}{zui4}[][HSK 6]
    \definition{s.}{crime; culpa | falha; culpa | sofrimento; dor; dificuldade | má conduta; transgressão; negligência | agonia; dor; sofrimento}
    \definition{v.}{colocar a culpa em alguém; culpar}
  \end{Phonetics}
\end{Entry}

\begin{Entry}{罪人}{13,2}{⽹,⼈}
  \begin{Phonetics}{罪人}{zui4ren2}
    \definition{s.}{culpado; infrator; pecador}
  \antonymref{功臣}{gong1chen2}
  \end{Phonetics}
\end{Entry}

\begin{Entry}{罪犯}{13,5}{⽹,⽝}
  \begin{Phonetics}{罪犯}{zui4fan4}[][HSK 7-9]
    \definition[名,个]{s.}{culpado; criminoso; infrator; pessoas que cometem crime}
  \synonymref{囚犯}{qiu2fan4}
  \synonymref{罪人}{zui4ren2}
  \antonymref{功臣}{gong1chen2}
  \end{Phonetics}
\end{Entry}

\begin{Entry}{罪行}{13,6}{⽹,⾏}
  \begin{Phonetics}{罪行}{zui4xing2}
    \definition{s.}{crime; culpa; delito; atos criminosos}
  \end{Phonetics}
\end{Entry}

\begin{Entry}{罪恶}{13,10}{⽹,⼼}
  \begin{Phonetics}{罪恶}{zui4'e4}[][HSK 6]
    \definition{s.}{pecado; mal; crime; comportamento criminoso grave}
  \end{Phonetics}
\end{Entry}

\begin{Entry}{罪魁祸首}{13,13,11,9}{⽹,⿁,⽰,⾸}
  \begin{Phonetics}{罪魁祸首}{zui4kui2-huo4shou3}[][HSK 7-9]
    \definition[个,位,名]{expr.}{principal culpado (ou criminoso); arquicriminoso; líder de quadrilha criminosa, principal infrator (expressão idiomática); Figurativo: causa principal de um desastre}
  \end{Phonetics}
\end{Entry}

%%%%%%%%%% 置 %%%%%%%%%%
\subsection*{置}\addcontentsline{loh}{figure}{置}

\begin{Entry}{置}{13}{⽹}
  \begin{Phonetics}{置}{zhi4}[][HSK 7-9]
    \definition{v.}{colocar; pôr | configurar; estabelecer; instalar; organizar | comprar; adquirir}
  \synonymref{放}{fang4}
  \synonymref{搁}{ge1}
  \end{Phonetics}
\end{Entry}

\begin{Entry}{置疑}{13,14}{⽹,⽦}
  \begin{Phonetics}{置疑}{zhi4yi2}
    \definition{v.}{(geralmente na negativa) duvidar}
  \synonymref{怀疑}{huai2yi2}
  \synonymref{可谓}{ke3wei4}
  \end{Phonetics}
\end{Entry}

%%%%% EOF %%%%%

